\documentclass[11pt,a4paper]{article}

\usepackage[a4paper,margin=2.2cm]{geometry}
\usepackage[utf8]{inputenc}
\usepackage[T1]{fontenc}
\usepackage{booktabs}
\usepackage{xcolor}
\usepackage{hyperref}
\usepackage{amsmath}

\hypersetup{
    colorlinks=true,
    linkcolor=black,
    urlcolor=blue
}

\setlength{\parindent}{0pt}
\setlength{\parskip}{0.8em}
\setlength{\emergencystretch}{2em}

\begin{document}

\begin{center}
    {\Large Informationsvisualisierung -- Übung 10: Implizite Baumdarstellungen (Treemaps)}
\end{center}

\section*{Lauffähiger Code}

\href{https://github.com/yaniktfl/iv-repo}{https://github.com/yaniktfl/iv-repo}

\section*{Ausführung}

\begin{verbatim}
cd 10
elm reactor
\end{verbatim}

Anschließend im Browser \texttt{http://localhost:8000} öffnen und dort
\texttt{Main.elm} auswählen.

\section*{Datenstruktur und Laden}

Als Datenstruktur wird wie vorgegeben die Baum-Implementierung
\texttt{zwilias/elm-rosetree} (\texttt{Tree}) genutzt. Jeder Knoten trägt ein
Paar \texttt{( String, Int )} aus Name und Wert. Die Datei \texttt{flare.json}
wird beim Start per \texttt{Http.get} geladen; ein rekursiver Decoder überführt
das geschachtelte JSON in einen \texttt{Tree ( String, Maybe Int )}:

\begin{verbatim}
treeDecoder : Decoder (Tree ( String, Maybe Int ))
treeDecoder =
    Decode.map3
        (\name value children ->
            Tree.tree ( name, value ) (Maybe.withDefault [] children))
        (Decode.field "name" Decode.string)
        (Decode.maybe (Decode.field "value" Decode.int))
        (Decode.maybe
            (Decode.field "children"
                (Decode.list (Decode.lazy (\_ -> treeDecoder)))))
\end{verbatim}

\texttt{Decode.lazy} ist nötig, da der Decoder sich für die Kinder selbst
aufruft. Nach dem Laden wird der Wert der Blätter über \texttt{computeSums}
rekursiv in die inneren Knoten aufaddiert, sodass jeder innere Knoten die
Summe seiner Kinder als Fläche kennt:

\begin{verbatim}
computeSums : Tree ( String, Int ) -> Tree ( String, Int )
computeSums tree =
    let
        children = Tree.children tree |> List.map computeSums
        ownValue = Tree.label tree |> Tuple.second
        childSum =
            children
                |> List.map (Tree.label >> Tuple.second)
                |> List.sum
    in
    Tree.tree ( Tree.label tree |> Tuple.first, ownValue + childSum ) children
\end{verbatim}

\section*{Aufgabe 10.1: Einfache Treemap (selbst implementiert)}

Die einfache Treemap wird vollständig selbst gezeichnet. Die Achsen werden
abwechselnd als Split-Achse genutzt, beginnend mit der x-Achse. Um die
Aufteilung der Kinder eines Knotens zu berechnen, wird der verfügbare Platz
entlang der Split-Achse linear auf die Summe der Kinder-Werte skaliert. Jedes
Kind erhält dadurch einen Anteil \texttt{childValue / childTotal} der Länge,
wird als Rechteck eingezeichnet und \texttt{drawTreeNode} ruft sich für dessen
Kinder rekursiv mit umgekehrter Split-Achse (\texttt{not splitX}) auf:

\begin{verbatim}
drawTreeNode : Bool -> Float -> Float -> Float -> Float
            -> Tree ( String, Int ) -> List (Svg Msg)
drawTreeNode splitX x y width height tree =
    ... children
        |> List.foldl
            (\child ( offset, acc ) ->
                let
                    childSize  = childValue / childTotal
                    nextWidth  = if splitX then width * childSize else width
                    nextHeight = if splitX then height else height * childSize
                    nextX      = if splitX then offset else x
                    nextY      = if splitX then y else offset
                    nextOffset =
                        if splitX then offset + nextWidth
                        else offset + nextHeight
                in
                ( nextOffset
                , acc ++ drawTreeNode (not splitX)
                            nextX nextY nextWidth nextHeight child ))
            ( if splitX then x else y, [] )
\end{verbatim}

Über die \texttt{foldl}-Akkumulation wandert der \texttt{offset} entlang der
aktuellen Split-Achse weiter, sodass die Kinder lückenlos nebeneinander (bzw.
übereinander) liegen. Die Treemap wird zunächst für den kleinen
\texttt{testTree} und anschließend für die Flare-Daten gezeichnet.

\section*{Aufgabe 10.2: Squarified Treemap (Hierarchy.treemap)}

Für die Squarified Treemap wird die Funktion \texttt{Hierarchy.treemap} aus dem
Modul \texttt{Hierarchy} der Bibliothek \texttt{gampleman/elm-visualization}
genutzt. Über \texttt{Hierarchy.tile Hierarchy.squarify} entstehen möglichst
quadratische Kacheln; \texttt{Hierarchy.size} legt die Zeichenfläche fest und
der Zugriffsfunktion \texttt{(Tuple.second >> toFloat)} entnimmt sie den
Flächenwert je Knoten:

\begin{verbatim}
positionedTree =
    Hierarchy.treemap
        [ Hierarchy.size canvasWidth canvasHeight
        , Hierarchy.paddingInner (\_ -> 1)
        , Hierarchy.paddingOuter (\_ -> 0)
        , Hierarchy.tile Hierarchy.squarify
        ]
        (Tuple.second >> toFloat)
        tree
\end{verbatim}

Der so positionierte Baum (\texttt{Tree PositionedNode} mit
\texttt{x}/\texttt{y}/\texttt{width}/\texttt{height}) wird mit
\texttt{Tree.depthFirstFold} in SVG-Rechtecke übersetzt. Labels werden nur
gezeichnet, wenn die Kachel groß genug ist (\texttt{labelIfRoom}).

\subsection*{Färbung ab Ebene 1 und ab Ebene 2}

Im Beispiel sind Teilbäume ab einer Ebene unter der Wurzel verschieden
gefärbt. Zusätzlich wird gefordert, Teilbäume ab der \emph{zweiten} Ebene unter
der Wurzel verschieden zu färben. Beide Modi werden über einen Typ
\texttt{ColorMode} getrennt. Die Farbe ergibt sich aus dem Pfad-Element auf der
gewählten Ebene: Bei \texttt{ColorByLevel1} entscheidet das Element an Index~1
(direktes Kind der Wurzel), bei \texttt{ColorByLevel2} das an Index~2 -- alle
Knoten innerhalb desselben Teilbaums teilen sich dadurch eine Farbe:

\begin{verbatim}
type ColorMode = ColorByLevel1 | ColorByLevel2

colorClassForNode mode ancestors node =
    let
        path = List.reverse (node :: ancestors)
                   |> List.map (.node >> Tuple.first)
        key =
            case mode of
                ColorByLevel1 ->
                    itemAt 1 path |> Maybe.withDefault "root"
                ColorByLevel2 ->
                    itemAt 2 path
                        |> Maybe.withDefault
                            (itemAt 1 path |> Maybe.withDefault "root")
    in
    "color-" ++ String.fromInt (paletteIndex key)
\end{verbatim}

Der \texttt{key} wird über \texttt{hashString} stabil auf einen von zwölf
Palettenplätzen (\texttt{modBy 12}) abgebildet, sodass gleiche Teilbäume
reproduzierbar dieselbe Farbe erhalten. Bei \texttt{ColorByLevel2} fällt ein
Knoten, der noch keine zweite Ebene besitzt, auf seine Ebene-1-Farbe zurück.

\section*{Überblick}

\begin{center}
\begin{tabular}{lll}
\toprule
Aufgabe & Verfahren & Layout \\
\midrule
10.1 einfache Treemap & selbst, x/y alternierend & lineare Skalierung \\
10.2 squarified Treemap & \texttt{Hierarchy.treemap} & \texttt{squarify}, Farben Ebene 1 \& 2 \\
\bottomrule
\end{tabular}
\end{center}

\end{document}
